\subsection{Actual standard-affinoid neighbourhoods}\label{subsec:fc-line-affinoid-basis}

\begin{lemma}[Finite closed linear-coordinate boxes]\label{lem:fc-finite-linear-box}
\lean{EnvelopeProperty.Berkovich.RelativeSpectrum.fc_isInducing_linearCoordinates, EnvelopeProperty.Berkovich.RelativeSpectrum.fc_exists_finite_linear_box}
\leanok
\uses{lem:fc-linear-convergence}
Let $x$ be an actual point of the analytic line and let $N$ be a neighbourhood
of $x$. There are finitely many centres $a_i$ and bounds $0\le l_i<u_i$
such that $|T-a_i|_x<u_i$, and either $l_i=0$ or $l_i<|T-a_i|_x$, and
\[
 \{y: l_i\le |T-a_i|_y\le u_i\text{ for every }i\}\subseteq N.
\]
There is at least one upper bound.
\end{lemma}

\begin{proof}
The map sending a point to all its linear-coordinate values induces the
analytic topology: convergence to any point is equivalent to convergence of
every linear value. A product neighbourhood therefore supplies finitely many
coordinate neighbourhoods. In each coordinate choose a sufficiently small
closed interval about its value, with a strictly larger upper endpoint and a
strictly smaller lower endpoint unless the value is zero. Those intervals lie
in the chosen coordinate neighbourhoods. Add the harmless upper bound
$|T|\le |T|_x+1$ if necessary to make the finite family nonempty.
\end{proof}

\begin{lemma}[Pruning finitely many open holes]\label{lem:fc-finite-hole-pruning}
\lean{EnvelopeProperty.Berkovich.RelativeSpectrum.fc_openDiscs_laminar, EnvelopeProperty.Berkovich.RelativeSpectrum.fc_disjoint_maximal_subfamily}
\leanok
\uses{lem:fc-inflation-linear}
A finite family of positive-radius open discs contained in a fixed closed
disc has a pairwise disjoint subfamily with the same union.
\end{lemma}

\begin{proof}
Two intersecting open discs are nested. Indeed, order their radii $r\le s$
and choose a point $z$ in their intersection. The constant-distance inequality
gives $|a-b|\le\max\{|T-a|_z,|T-b|_z\}<s$, and then every point of the
first disc belongs to the second by the ultrametric inequality. Regard the
finite family as a finite set of subsets, so repeated discs occur only once.
Retain its maximal members under inclusion. Every original member is contained
in a maximal member, so the union is preserved. Distinct maximal members are
disjoint by nesting.
\end{proof}

\begin{theorem}[Actual standard-affinoid neighbourhood basis]\label{thm:fc-standard-affinoid-basis}
\lean{EnvelopeProperty.Berkovich.RelativeSpectrum.fc_exists_standardAffinoid_for_finite_lower_bounds, EnvelopeProperty.Berkovich.RelativeSpectrum.fc_exists_standardAffinoid_nhds, EnvelopeProperty.Berkovich.RelativeSpectrum.fc_nhds_hasBasis_standardAffinoid}
\leanok
\uses{lem:fc-finite-linear-box, lem:fc-finite-hole-pruning, lem:fc-standard-boundary}
For every actual point $x$ and every neighbourhood $N$ of $x$, there is an
actual standard affinoid $U$ with $x\in\operatorname{int}(U)$ and $U\subseteq N$.
Consequently the interiors of actual standard affinoids form a neighbourhood
basis at every point.
\end{theorem}

\begin{proof}
Take the finite linear-coordinate box above and choose an upper disc of least
radius, say $D(a_0,r_0)$. It is contained in every other upper disc: their
centres have distance at most the maximum of their values at $x$, strictly
smaller than the larger upper radius, and $r_0$ is the least upper radius.
For each positive lower bound remove the open disc $D^-(a_i,l_i)$.
Discard those lower discs whose centres are outside $D(a_0,r_0)$. They do not
meet the outer disc: throughout the outer disc their linear value equals the
constant centre distance, which is also their value at $x$ and is greater
than $l_i$. Every retained centre belongs to the outer disc. Its lower radius
is strictly smaller than its value at $x$, which is at most $r_0$; hence its
open disc is contained in the outer disc. Prune the retained holes by the
preceding lemma, and enumerate the resulting finite set of discs. Their actual
centres and positive radii, together with the original outer disc, give the
standard-affinoid data and all its required inequalities and disjointness.
Its carrier is the outer disc with exactly the original relevant lower discs
removed, so it lies in the finite box and hence in $N$.
At $x$ every upper inequality and every retained lower inequality is strict.
The corresponding finite intersection of strict coordinate inequalities is an
open neighbourhood of $x$ contained in the carrier. Thus $x$ is interior.
The same construction for every $N$ gives the asserted neighbourhood basis.
\end{proof}
